\section{Đảo câu nguồn}
\label{sec:ch05-dao-cau-nguon}

\index{đảo câu nguồn!phát biểu}%
Trong ba chỗ mà mô hình thật của bài khác mô hình trong mục 2, hai chỗ đầu là
chuyện kích thước. Chỗ thứ ba thì không đụng gì tới kiến trúc:
\tn{đảo câu nguồn}{source reversal}, tức đảo ngược thứ tự từ của
\tn{câu nguồn}{source sentence}, và chỉ câu nguồn, trước khi đưa vào mạng. Thay
vì học ánh xạ \(a, b, c\) sang \(\alpha, \beta, \gamma\), mạng được hỏi ánh xạ
\(c, b, a\) sang \(\alpha, \beta, \gamma\).

Bài gọi đây là một trong những đóng góp kỹ thuật chính của nó, và con số họ đưa
ra không nhỏ: perplexity trên tập test giảm từ 5.8 xuống 4.7, BLEU tăng từ 25.9
lên 30.6.

\index{đảo câu nguồn!lập luận theo khoảng cách}%
Lập luận thì đáng đọc kỹ hơn con số, vì nó là lập luận về khoảng cách và nó
kiểm được. Nối câu nguồn với \tn{câu đích}{target sentence} lại thành một
chuỗi, rồi hỏi mỗi từ nguồn
cách từ đích tương ứng của nó bao xa. Bài viết rằng đảo câu nguồn để nguyên
khoảng cách \emph{trung bình} giữa các cặp từ tương ứng, rồi nói tiếp:
\enquote{however, the first
few words in the source language are now very close to the first few words in
the target language, so the problem's minimal time lag is greatly reduced}.

Chữ \enquote{minimal time lag} ấy đã có trong chương~\ref{ch:hien-tuong}, ở hộp
cầu nối về Hochreiter 1991. Nó là quãng ngắn nhất mà tín hiệu lỗi phải đi được
để mạng bắt đầu học được cái gì đó. Encoder đọc xong thì dừng ở từ \emph{cuối}
của câu nguồn, và decoder viết ra từ \emph{đầu} của câu đích ngay sau đó. Không
đảo thì từ đích đầu tiên và từ nguồn ứng với nó nằm ở hai đầu đối diện của chuỗi
ghép, tức xa nhau nhất có thể. Đảo thì chúng thành cạnh nhau.

\subsection*{Lập luận ấy dự đoán một thứ sắc hơn \enquote{đảo thì tốt hơn}}

Nếu tác dụng của phép đảo đến từ khoảng cách, thì phần lợi phải rơi vào câu dài.
Với câu ngắn thì mọi từ đã đủ gần rồi, đảo hay không cũng thế, nên ở đó phép đảo
phải gần như không làm gì.

Đó là một dự đoán, không phải một cách diễn đạt lại. Nó sai được. Bảng dưới đây
chạy ba seed mỗi vế, và in ra từng seed chứ không chỉ trung bình.

\begin{measured}{đảo câu nguồn, ba seed mỗi vế}
\begin{minted}{text}
source     seed  loss     exact    short    long
raw        0     0.0644   0.6733   1.0000   0.3553
raw        1     0.1392   0.5367   1.0000   0.0855
raw        2     0.0997   0.5933   1.0000   0.1974
reversed   0     0.0251   0.8533   1.0000   0.7105
reversed   1     0.0652   0.6433   0.9797   0.3158
reversed   2     0.0353   0.8433   1.0000   0.6908

mean over seeds:
source     exact    short    long
raw        0.6011   1.0000   0.2127
reversed   0.7800   0.9932   0.5724

reversed minus raw:
  exact  +0.1789
  short  -0.0068
  long   +0.3596
\end{minted}

Đo bằng \code{experiments/ch05\_reverse.py} tại tag \code{ch05}.
\end{measured}

\index{đảo câu nguồn!phần lợi rơi vào câu dài}%
Dự đoán đúng, và đúng gọn hơn tôi tưởng. Câu ngắn đứng yên: 1.0000 thành 0.9932,
tức lệch đi một câu trong một trăm bốn tám, và lệch theo chiều xấu đi. Câu dài
đi từ 0.2127 lên 0.5724. Toàn bộ 0.1789 điểm mà phép đảo mang lại đến từ nhóm
câu dài, đúng nhóm mà lập luận khoảng cách nói nó phải đến.

\subsection*{Đọc bảng này theo cặp, không theo trung bình}

\index{đo lường!so theo cặp cùng seed}%
Ba dòng \code{reversed} tản từ 0.6433 tới 0.8533. Ba dòng \code{raw} tản từ
0.5367 tới 0.6733. Hai dải ấy chồng lên nhau: seed 1 có đảo (0.6433) còn thấp
hơn seed 0 không đảo (0.6733). Nếu tôi chỉ chạy một seed mỗi vế và rút phải hai
seed đó, tôi đã kết luận phép đảo làm hại.

Cách đọc đúng là so theo cặp cùng seed, vì hai vế cùng seed dùng chung cả khởi
tạo lẫn thứ tự minibatch: 0.8533 so 0.6733, rồi 0.6433 so 0.5367, rồi 0.8433 so
0.5933. Ba trên ba, và trên cột \code{long} cũng ba trên ba. Đó là bằng chứng,
còn trung bình một mình thì không.

Chuyện này không phải mẹo thống kê mà là hệ quả của quy mô. Ở đây một lần huấn
luyện là một con số, không phải một xu hướng, và chương~\ref{ch:lstm} đã học bài
ấy một lần rồi ở phép đo cosine của nó.

\subsection*{Chỗ bài tự nhận là bất ngờ}

\index{đảo câu nguồn!điều bài không dự đoán được}%
Có một đoạn trong mục 3.3 mà tôi thích, vì nó là chỗ hiếm hoi một bài báo ghi
lại dự đoán sai của chính mình. Các tác giả viết rằng ban đầu họ tin phép đảo
chỉ làm mạng tự tin hơn ở \emph{đầu} câu đích và kém tự tin hơn ở
\emph{cuối} câu, tức một phép đổi chỗ chứ không phải một phép cải thiện. Rồi họ
đo và thấy mạng huấn luyện trên câu đảo làm tốt hơn hẳn trên câu dài, và họ kết
luận rằng phép đảo cho ra những mạng dùng bộ nhớ tốt hơn
(\enquote{LSTMs with better memory utilization}).

Bảng trên nói cùng một chuyện bằng một phép đo khác: cột \code{short} không mất
gì, cột \code{long} được rất nhiều. Không có chỗ nào tệ đi để đổi lấy chỗ tốt
lên. Nếu đây chỉ là đổi chỗ, phải có một cột đi xuống, và không có.

Còn một câu nữa trong mục 5 đáng ghi, vì nó là một chỗ bài để ngỏ: họ không
huấn luyện nổi một mạng hồi quy thường trên bài toán không đảo, nhưng tin rằng
một mạng như thế sẽ huấn luyện được dễ dàng khi câu nguồn đã đảo. Nguyên văn:
\enquote{we believe that a standard RNN should be easily trainable when the
source sentences are reversed}. Và họ nói thẳng ra là không kiểm bằng thí
nghiệm. Bài tập cuối chương có một câu hỏi về chỗ đó.
